%%
%% This is file `mathfont_example_kelvinch.tex',
%% generated with the docstrip utility.
%%
%% The original source files were:
%%
%% mathfont_code.dtx  (with options: `kelvinch')
%% 
%% This file is from version 2.0 of the free and open-source
%% LaTeX package "mathfont," to be used with the XeTeX or
%% LuaTeX engines. (As of version 2.0, LuaTeX is recommended.)
%% 
%% Copyright 2018-2021 by Conrad Kosowsky
%% 
%% This file may be distributed and modified under the terms
%% of the LaTeX Public Project License, version 1.3c or any
%% later version. The most recent version of this license is
%% available online at
%% 
%%           https://www.latex-project.org/lppl/.
%% 
%% This Work has the LPPL status "maintained," and the current
%% maintainer is the package author, Conrad Kosowsky. He can
%% be reached at kosowsky.latex@gmail.com.
%% 
%% THIS SOFTWARE IS PROVIDED "AS IS" WITHOUT ANY EXPRESS
%% OR IMPLIED WARRANTY, INCLUDING THE IMPLIED WARRANTIES
%% OF MERCHANTABILITY OR FITNESS FOR A PARTICULAR PURPOSE.
%% IF THE SOFTWARE IS DEFECTIVE, YOU AGREE TO ASSUME THE
%% COST FOR ANY REPAIR OR CORRECTION.
%% 
%% BY USING OR DISTRIBUTING THIS SOFTWARE, YOU AGREE
%% TO RELEASE THE PACKAGE AUTHOR FROM ANY LIABILITY FOR
%% DAMAGES ARISING OUT OF YOUR USE OF THE SOFTWARE. YOU
%% AGREE TO DO SO TO THE MAXIMUM EXTENT ALLOWED UNDER
%% APPLICABLE LAW AND EVEN IF THE PACKAGE AUTHOR HAS
%% BEEN ADVISED OF THE POSSIBILITY OF SUCH DAMAGES.
%% 
%% See also the "No Warranty" section of the LaTeX Project
%% Public License. In releasing the package author from
%% liability, you also release from liability any third
%% parties who distribute the software under the terms
%% of the LaTeX Project Public License. Derivative works
%% based on this package may come with their own license or
%% terms of use, and the package author is not responsible
%% for any third-party software.
%% 
%% For more information, see the mathfont_code.dtx.
%% 
%% Happy TeXing!
%% 
\documentclass[12pt]{article}
\makeatletter
\usepackage[margin=1in]{geometry}
\usepackage{innerscript}
\usepackage{multicol}
\usepackage{amsmath}
\DeclareMathOperator{\Res}{Res}
\usepackage{mathfont}
\mathfont[radical]{Crimson}
\setfont{Kelvinch}
\IntegralItalicFactor{500}
\SurdHorizontalFactor{900}
\parindent=0pt\relax
\begin{document}

\let\footnote\@gobble
\def\showabstract{0}
\let\textsf\relax
\let\ttfamily\relax
\def\documentname{Example---Kelvinch}

\input mathfont_heading.tex

This is Kelvinch with Crimson for the radical signs. Testing. Testing. Brown foxes quickly jump over dazzling does and harts. This document shows an example of mathfont in action. Unfortunately, there are many more equations in the world than I have space for here. Nevertheless, I hope I hit some of the highlights. Happy \TeX ing!

%%
%% This is file `mathfont_equations.tex',
%% generated with the docstrip utility.
%%
%% The original source files were:
%%
%% mathfont_code.dtx  (with options: `equations')
%% 
%% This file is from version 2.0 of the free and open-source
%% LaTeX package "mathfont," to be used with the XeTeX or
%% LuaTeX engines. (As of version 2.0, LuaTeX is recommended.)
%% 
%% Copyright 2018-2021 by Conrad Kosowsky
%% 
%% This file may be distributed and modified under the terms
%% of the LaTeX Public Project License, version 1.3c or any
%% later version. The most recent version of this license is
%% available online at
%% 
%%           https://www.latex-project.org/lppl/.
%% 
%% This Work has the LPPL status "maintained," and the current
%% maintainer is the package author, Conrad Kosowsky. He can
%% be reached at kosowsky.latex@gmail.com.
%% 
%% THIS SOFTWARE IS PROVIDED "AS IS" WITHOUT ANY EXPRESS
%% OR IMPLIED WARRANTY, INCLUDING THE IMPLIED WARRANTIES
%% OF MERCHANTABILITY OR FITNESS FOR A PARTICULAR PURPOSE.
%% IF THE SOFTWARE IS DEFECTIVE, YOU AGREE TO ASSUME THE
%% COST FOR ANY REPAIR OR CORRECTION.
%% 
%% BY USING OR DISTRIBUTING THIS SOFTWARE, YOU AGREE
%% TO RELEASE THE PACKAGE AUTHOR FROM ANY LIABILITY FOR
%% DAMAGES ARISING OUT OF YOUR USE OF THE SOFTWARE. YOU
%% AGREE TO DO SO TO THE MAXIMUM EXTENT ALLOWED UNDER
%% APPLICABLE LAW AND EVEN IF THE PACKAGE AUTHOR HAS
%% BEEN ADVISED OF THE POSSIBILITY OF SUCH DAMAGES.
%% 
%% See also the "No Warranty" section of the LaTeX Project
%% Public License. In releasing the package author from
%% liability, you also release from liability any third
%% parties who distribute the software under the terms
%% of the LaTeX Project Public License. Derivative works
%% based on this package may come with their own license or
%% terms of use, and the package author is not responsible
%% for any third-party software.
%% 
%% For more information, see the mathfont_code.dtx.
%% 
%% Happy TeXing!
%% 
\begin{multicols}{2}
\let\medskip\relax

Black-Scholes Equation
\[
\frac{\partial V}{\partial t}+\frac12\sigma^2S^2\frac{\partial^2V}{\partial S^2}=
  rV-rS\frac{\partial V}{\partial X}
\]
\medskip

Cardano's Formula/Cubic Formula
\begin{align*}
t_i&=\omega_i\sqrt[3]{-\frac q2+\sqrt{\frac{q^2}4+\frac{p^3}{27}}}\\
  &\qquad\qquad{}+\omega_i^2\sqrt[3]{-\frac q2-\sqrt{\frac{q^2}4+\frac{p^3}{27}}}
\end{align*}
\medskip

Einstein's Field Equation (General Relativity)
\[
R_{\mu\nu}-\frac12Rg_{\mu\nu}+\Lambda g_{\mu\nu}=\frac{8\pi G}{c^4}T_{\mu\nu}
\]
\medskip

First Isomorphism Theorem
\[
\phi(X)\cong X/\ker(\phi)
\]
\medskip

Gauss-Bonnet Formula
\[
\int_MK\ dA+\int_{\partial M}k_g\ ds=2\pi\chi(M)
\]
\medskip

Maxwell's Equations
\begin{align*}
\nabla\bullet\mathbf E&=\frac\rho{\epsilon_0}&
\nabla\bullet\mathbf B&=0\\
\nabla\times\mathbf E&=-\frac{\partial\mathbf B}{\partial t}&
\nabla\times\mathbf B&=\mu_0\left(\mathbf J+
  \epsilon_0\frac{\partial\mathbf E}{\partial t}\right)
\end{align*}
\medskip

Michaelis-Menten Model
\[
v=\frac{d[P]}{dt}=V\frac{[S]}{K_M+[S]}
\]

%% next column begins here

Navier-Stokes Equation
\begin{align*}
\frac{\partial}{\partial t}(\rho\mathbf{u})+\nabla\bullet
  (\rho \mathbf{u}\otimes \mathbf{u})={}
  &{-}\nabla\bar p+\mu\nabla^2\mathbf{u}\\
  &+\frac13\mu\nabla(\nabla\bullet u)+
  \rho\mathbf{g}
\end{align*}
\medskip

Quadratic Formula
\[
x=\frac{-b\pm\sqrt{b^2-4ac}}{2a}
\]
\medskip

Ramanujan's Approximation for $\Gamma$
\[
\Gamma(1+x)\approx\sqrt\pi\,x^xe^{-x}\,\sqrt[6]{8x^3+4x^2+x+\frac1{30}}
\]
\medskip

Residue Theorem
\[
\frac1{2i\pi}\int_{\gamma}f(z)\ dz=\sum_{k=1}^n\Res_{a_k}(f)
\]
\medskip

Riemann Zeta Function
\begin{align*}
\zeta(z)&=\sum_{i=1}^\infty\frac1{z^i}
  =\frac1{\Gamma(z)}\int_0^\infty\frac{x^{z-1}}{e^x-1}\ dx\\
  &=2^z\pi^{z-1}\sin\left(\frac{\pi z}2\right)\,\Gamma(1-z)\,\zeta(1-z)
\end{align*}
\medskip

Schrodinger Equation
\[
i\hbar\frac d{dt}|\Psi(t)\fakerangle=\hat H|\Psi(t)\fakerangle
\]
\medskip

Lorentz Transformation (Special Relativity)
\[
t'=\left(t-\frac{vx}{c^2}\right)\frac1{\sqrt{1-\frac{v^2}{c^2}}}
\]

\end{multicols}
\endinput
%%
%% End of file `mathfont_equations.tex'.


\end{document}
\endinput
%%
%% End of file `mathfont_example_kelvinch.tex'.
